% SOURCE docs/release/2026-10-29/source/docs/writing-repair/2026-10-29/investment.tex CHAR 0
\chapter{The first investment and the evidence that would reverse it}
\label{rewrite:investment-decision}

The immediate proposal is a minimal comparative integration for a fixed
offer study, with explicit population weights, information dates and
a complete financial report. Reuse a transparent estimator and retain
a documented alternative. The seventeen-area comparison establishes
which further capabilities are plausible candidates and why they need
separate assumptions. It is not a reason to build all seventeen before
evaluating the first study.

\section{Compare equivalent completed studies}
There are four economically serious alternatives. Expert composition
uses existing methods and bespoke study documentation. Minimal integration
adds a reusable request and result description around selected methods.
A broad shared service adds more methods, languages, maintenance and
operating commitments. A commercial route can provide supported
procedures and existing institutional integration, but its relevant
entitlements, actual prices and matched capabilities are unknown here.
The detailed commercial discussions preserve documented features and
unknowns; an unobserved feature is not evidence that it is absent.

The proposed first comparison selects minimal integration against expert
composition. Its potential gains come from repeated target specification,
information-date checks, error detection and reporting. Fitting a regression
may occupy little of the total work. Any proposed saving must therefore
be measured over preparation, correction, interpretation and failed
attempts. Bespoke causal reasoning remains part of every alternative.

Let $N$ be annual studies genuinely eligible for this restricted route;
$h$ hours saved per completed study at equivalent quality; $w$ loaded
hourly cost; $D$ reduction in expected decision loss per study; $C$ annual
operating cost; and $K$ annualised development and adoption cost. Then
\[
 B=N(wh+D)-C-K.
\]
The conditional break-even study volume is $(C+K)/(wh+D)$ only when
$wh+D>0$. A negative denominator cannot be repaired by doing more studies.
The illustrative investment account in Section~\ref{r3:investment}
retains its actual stipulated inputs and arithmetic, including the
26.74-study threshold. That threshold belongs to that account. It must
not be imported as an observed requirement for the newly selected scope.

The deposit example's \pounds3 population contribution concerns one
constructed customer action. It is not $D$, the improvement in decision
loss due to choosing one analytical workflow over another. Counting
that whole contribution as a software benefit would confuse policy value
with the incremental value of analysis. Similarly, correction hours
counted in $h$ must not be counted again in $D$. These distinctions
prevent an attractive worked example from becoming an invented ROI.



% SOURCE docs/release/2026-10-29/source/docs/writing-repair/2026-10-29/src/830f6276ba2598e7.tex CHAR 0
\section{The multi-year investment account and institutional variation}
\label{r3:investment}

A common causal product earns its cost by improving repeated scientific
work. The appropriate alternative is a capable quantitative team composing
the same open-source methods with ordinary engineering, documentation and
review. Such a team can produce good causal studies without a new platform.
The investment question is whether the shared representations and qualified
operations reduce the cost of equivalent work or enable additional valuable
decisions, after construction, support and adoption are included.

\subsection{Define equivalent output before pricing reuse}
\label{r3:equivalent-output}

An equivalent completed study has a reconstructed population and observation
clock, a clear intervention and financial outcome, a defensible identification
argument or useful nonidentification conclusion, appropriate inference,
reproducible artifacts and review. The comparison must include negative
results. An early finding that a requested target is unidentified can save
substantial work; forcing every study to return a numerical effect would
reward scientific error.

The proposed service would help analysts reconstruct what was known when
a decision was made, keep the population and outcome consistent, document
why an effect is identified, and carry uncertainty into the financial
comparison. It would reuse existing numerical methods where their
assumptions fit. Each business would still specify its products, costs and
constraints. A mortgage policy, an individual collections decision and a
macroeconomic shock need not share the same economic quantity of interest.

Comparator costs should include the same required outputs and review
quality. Reduced effort from omitting uncertainty or failing to investigate
weak support is not a saving. Conversely, building a generic model interface
for a one-off study can cost more than a small, clearly documented expert
analysis. The portfolio of genuinely repeated requests, not the total
number of database rows, determines the economic opportunity.

\subsection{An auditable cost model}
\label{r3:cost-model}

For alternative $j\in\{E,P\}$, expert composition or shared product,
let $K_j$ be initial construction cost, $F_j$ annual shared cost,
$c_j$ marginal cost per equivalent completed study, and $n_t$ completed
studies in year $t$. With horizon $H$ and discount rate $r$,
\begin{equation}
 C_j=K_j+\sum_{t=1}^{H}\frac{F_j+c_jn_t}{(1+r)^t}.
 \label{r3:present-cost}
\end{equation}
The categories are disjoint: method maintenance and common infrastructure
belong in $F_j$, while study-specific preparation, fitting, review and
rework belong in $c_j$. Failed-study costs must be included in one
category; a completed-study denominator cannot make them disappear.

Use the following illustrative USD-million assumptions to expose the
break-even logic:
\[
 (K_E,F_E,c_E)=(1,.6,.15),\qquad
 (K_P,F_P,c_P)=(4,1.8,.06),\qquad H=3,\quad r=.10.
\]
These are planning inputs, not observed bank prices, measured productivity
or enterprise vendor quotations. Let
$a_H=\sum_{t=1}^H(1+r)^{-t}$. At constant annual volume $n$, the cost-only
threshold, when $c_E>c_P$, is
\begin{equation}
 n^*=\frac{(K_P-K_E)/a_H+(F_P-F_E)}{c_E-c_P}.
 \label{r3:reuse-threshold}
\end{equation}
If the marginal saving is zero or negative, this particular reuse argument
does not establish a positive-volume crossover. Under the stipulated
numbers, $a_H=2.48685$ and $n^*=26.74$ studies annually.

\begin{table}[htbp]\centering\small


% SOURCE docs/release/2026-10-29/source/docs/writing-repair/2026-10-29/src/3a8f90066fa1af4f.tex CHAR 0
\begin{tabular}{@{}lrrr@{}}
\toprule Scenario & Expert cost & Platform cost & Platform saving\\\midrule
Low reuse & 5.476 & 9.670 & -4.194\\
Moderate reuse & 11.445 & 12.057 & -0.613\\
High reuse & 15.921 & 13.848 & 2.073\\
Adoption ramp & 8.402 & 10.840 & -2.438\\
\bottomrule\end{tabular}


% SOURCE docs/release/2026-10-29/source/docs/writing-repair/2026-10-29/src/830f6276ba2598e7.tex CHAR 3623

\caption{Three-year present cost in illustrative USD millions. Annual
volumes are 8, 24 and 36, or the ramp $(4,16,30)$. Both alternatives
deliver the same specified scientific outputs.}
\label{r3:reuse-economics}
\end{table}

At eight annual studies the product costs about USD\,4.19 million more.
At 24 it still costs about USD\,.61 million more. At 36 it saves about
USD\,2.07 million. The ramp leaves a USD\,2.44 million premium even
though final-year volume exceeds the constant-volume threshold:
construction occurs before most reuse. Halving marginal savings doubles
the threshold to 53.47 studies; increasing the annual shared-cost
premium to USD\,2 million raises it to 35.63.

\begin{figure}[htbp]\centering
\includegraphics[width=.88\textwidth]{derivations/reuse_economics.pdf}
\caption{The crossover follows from the declared account. It is not
a forecast of demand or evidence that the shared layer already saves effort.}
\end{figure}

The simple model holds study mix fixed. A more realistic portfolio has
classes $k$ with costs $c_{jk}$ and volumes $n_{tk}$, replacing
$c_jn_t$ by $\sum_k c_{jk}n_{tk}$. A few expensive bespoke latent
models may have little shared work with many routine rollouts.
Reuse should therefore be measured by operation and data contract, not
assumed from a common causal-inference label.

\subsection{Separate business value from construction savings}
\label{r3:incremental-value}

Let $D_t$ be incremental realized decision value attributable to the
product relative to the equivalent expert workflow. Then
\begin{equation}
 \Delta\operatorname{NPV}
 =C_E-C_P+\sum_{t=1}^{H}\frac{D_t}{(1+r)^t}.
 \label{r3:incremental-npv}
\end{equation}
The first term already counts workflow savings. Released analyst time
must not be added again as business value. If it permits extra studies,
their costs, capacity use, outcome lags and incremental decisions
belong in the expanded account.

At 24 annual studies the illustrative cost premium requires about
USD\,.246 million of additional constant annual decision value.
The adoption ramp requires about USD\,.980 million. The size of a
bank's balance sheet does not establish either amount. The expert
comparator may already identify the same profitable policy or reject
the same loss-making offer.

The deposit study illustrates the attribution problem. Finding a
negative contribution can prevent a loss only if the finding changes
what would otherwise have been done. The prevented loss cannot be
credited to the product merely because its report contains the estimate.
Likewise the simulated collections gain is a property of the declared
environment; it supplies no observed revenue to this investment account.
Latent-state research can be worthwhile even while operational value
remains uncertain.

Costs and gains also mature on different clocks. Data access,
implementation and review are paid before customer outcomes are
observed. A six-month retention measure may mature during an early
evaluation; credit losses or a persistent pricing response may not.
A continuation model supplies a conditional forecast for those outcomes,
not an observation. Report the cash timing and model dependence explicitly.

\subsection{Measure the common-product hypothesis}
\label{r3:adoption-experiment}

The complete-workflow measurements in Section~\ref{s4:workflow-results}
refine the resource account but do not estimate its labor terms. A response
workflow over 3.85 million versioned records takes about 3.38 seconds on the
recorded host; a 3,200-customer latent workflow with 19 bootstrap refits takes
about 48.27 seconds. Both the product and expert-composition arms can use
these same numerical operations. Pricing one arm from a bare estimator call
and the other from reconstruction plus inference would be an unequal-output
comparison. The illustrative per-study dollar inputs above therefore remain
unchanged pending matched measurements of preparation, review and rework.

The first tranche should compare matched tasks in the intervention,
sequential-decision and dynamic-systems workspaces. Agree their outputs,
scientific scope and reviewer rubric before measuring effort. Give the
expert-composition arm the same sources, data access and expertise, and
allow sensible reusable templates. Counterbalance task order where
possible to reduce learning effects.

Measure analyst and reviewer hours, calendar delay, machine cost,
failed starts, rework and unresolved material findings separately.
Ask a second analyst to reconstruct the population, target and financial
conclusion from the exported artifacts. Deliberately introduced clock,
weight and sign errors can test whether review finds the defects.
This tests the proposed product's specific value proposition, rather
than whether a polished interface is pleasant.

A planning target of 30\% lower median analyst/reviewer effort can be
useful if review quality and output scope are held fixed. It remains
a hypothesis until measured across repeated uses and teams. A successful
first-author demonstration is insufficient evidence of institutional
reuse. Conversely, more early rejections can be a benefit if they
correctly avoid expensive unidentified studies.

Progression can be written as
\[
 n_t^{\rm usable}=n_t^{\rm started}
 p_t^{\rm data}p_t^{\rm design}p_t^{\rm execution}p_t^{\rm review},
\]
where each probability conditions on reaching the preceding stage.
This is the chain rule, not an independence assumption. It separates
volume aspirations from usable output and makes adoption bottlenecks
visible. Charges for unsuccessful studies remain in the cost account.

\subsection{A staged construction budget and stopping rules}
\label{r3:investment-sequence}

An illustrative first tranche assigns six people for six months:
one product/economics lead, two causal/econometric scientists, two
data/platform engineers and one application engineer. At USD\,25,000
per loaded person-month, staff cost is USD\,900,000; USD\,50,000
compute/tooling and 20\% contingency give USD\,1.14 million.
This tranche is included in the illustrative USD\,4 million initial
product cost, not added to it. Data-owner, independent validation,
security and application-owner participation are unpriced and must
be estimated consistently for both alternatives before a bank funds work.

The tranche should implement the shared observation/request/evidence
objects and qualify a narrow adjusted/cohort and observed-history
policy path. Controlled-state fitting and temporal discovery should
enter the architecture immediately through finite, explicitly conditional
research operations. Six people cannot qualify every method in the
44-project survey for production in six months.

The next tranche depends on concrete outcomes. Continue if repeated
requests share observation and scientific contracts, reviewers can
reconstruct results, the required estimators/inference pass their
acceptance exercises, and measured effort supports the economic case.
Narrow the scope if most work remains application-specific. Retain
expert composition when reuse is low. Stop a particular policy analysis
when its data cannot distinguish materially different decisions and
additional evidence is not worth its cost.

These are different decisions: discontinuing an unidentified analysis,
deferring a numerical adapter, and declining investment in the shared
product should not be collapsed into a single project-failure flag.
The source archive and scientific representations remain useful even
if one deployment case is rejected.

\subsection{An independent offering must price institutional variation}
\label{r3:commercial-product}

An independent product can distribute the scientific core as an embedded
library and provide the three workspaces with deployment, support and
upgrade services. A design-partner engagement is a plausible first
commercial form: a named study portfolio, paid integration scope,
reproducible evidence export and a recurring support boundary.
No willingness-to-pay, sales-cycle or customer-retention estimate is
established by this proposal.

Each institution adds semantic mapping, legal-entity structure, observation
systems, scientific review and operational requirements. Code reuse
reduces some common work but does not remove those costs. A commercial
comparison should price onboarding, customer-specific adapters, method
maintenance, support and validation participation, and state which
customizations remain owned by the customer.

Procurement should compare products on the same reconstructed studies,
supported identification routes, uncertainty objects, data-locality
requirements and export/replay obligations. Repository activity or an
enterprise demonstration is not evidence that these scientific
operations are available. The source review establishes open-source
construction alternatives; it does not assert uninspected private-vendor
capabilities or prices.

The multi-year account prices the reuse hypothesis rather than a causal
method in isolation. Its cost and volume assumptions are sensitivity inputs
to the decision in this chapter, not observed bank demand.


% SOURCE docs/release/2026-10-29/source/docs/writing-repair/2026-10-29/src/014c93f1736bc825.tex CHAR 0
\hypertarget{pd15:investment}{%
\section{The value of obtaining additional evidence}\label{pd15:investment}}

\hypertarget{proposal-acceptance-and-later-qualification}{%
\subsection{The evidence needed before investing}\label{proposal-acceptance-and-later-qualification}}

The investment decision requires a comparison of feasible alternatives,
their costs and the evidence for their expected benefits. Reproducible
public and simulated studies can support that comparison within their
stated assumptions. A restricted result may justify a narrow prototype
while leaving broader research questions unresolved.

Public evidence informs three choices: which existing methods to reuse,
where language assistance is useful, and how analysts should conduct a
reproducible study. A dataset can reveal a timing error or support a
controlled comparison without settling every question about those choices.

Later bank qualification asks a different question: do the assumptions,
precision, workflows and economics hold for a specific institution? That
stage may need genuine offer histories, logging information, independent
labels, owner-reviewed semantics, production constraints and prospective
outcomes. The relevant input should be requested only for the capability
and target population under test. An internal bank pilot is not required
to demonstrate that the proposal has a reproducible public baseline and
a sensible experiment design.

\hypertarget{a-value-of-information-decision}{%
\subsection{A value-of-information
decision}\label{a-value-of-information-decision}}

Let \(d\in\{\text{stop},\text{build},\text{pilot}\}\) denote an
investment choice, \(\theta\) unresolved technical and economic
conditions, and \(U(d,\theta)\) its net value. Let \(E\) denote current
public/prototype evidence and \(Z\) the possible result of a proposed
experiment. A planning expression for the expected value of sample
information is

\[
\operatorname{EVSI}
=\mathbb E_{Z\mid E}\!\left[\max_d\mathbb E\{U(d,\theta)\mid E,Z\}\right]
-\max_d\mathbb E\{U(d,\theta)\mid E\}.
\]

An experiment merits consideration when its expected decision
improvement exceeds its full incremental cost, including data
preparation, customer burden where relevant, security review, labeling,
analyst time and opportunity cost. This is a project decision framework,
not an empirically estimated ROI. Where a probability model for
\(\theta\) is unjustified, use a scenario table and sensitivity bounds
instead of a precise-looking EVSI number.

A simple illustrative calculation shows the sequencing. Suppose, purely
for planning, a small bank study costing 20,000 could avert an
unsuitable 200,000 build in 20\% of scenarios, with no other benefits or
costs. Its gross avoided-spend value would be 40,000 and its net would
be 20,000 under those assumptions. If the study changes the build
decision in only 2\% of scenarios, the same gross value is 4,000 and it
is not justified on that ground. These are stipulated arithmetic
examples, not quotations, bank probabilities or forecasts.

The decision should identify the uncertainty that matters. If a public
benchmark shows repeated action-clock mistakes, improve the formulation
layer before requesting confidential data. If a correct estimator is too
slow on modest data, fix or replace it before commissioning a platform
pilot. If all plausible effects lead to the same action, a costly
precision improvement may have little value. If policy value reverses
when a missing hardship variable is plausible, a carefully scoped bank
measurement study can have high value even when the public-data
implementation is technically sound.

\hypertarget{conditions-for-moving-to-expensive-evidence}{%
\subsection{Conditions for moving to expensive
evidence}\label{conditions-for-moving-to-expensive-evidence}}

An actionable next-experiment recommendation specifies the proposed
intervention, target population, candidate baseline, decision-loss
threshold and the smallest information addition that could alter the
decision. It also states an explicit stop or redesign outcome. A bank
experiment should not be recommended merely to fill a missing-evidence
cell.

For funding, the missing addition may be time-stamped offers and an
assignment mechanism, not more aggregate deposits. For servicing, it may
be feasible-action and hardship information at each decision time. For
discovery, it may be one feasible intervention that distinguishes
competing graphs. For observation construction, it may be a
representative independently resolved sample of uncertain identities
with known inclusion probabilities. Public data help design each
measurement and estimate its computational burden, while the bank
supplies the later population-specific truth.

Authentic institutional approval, independent acceptance, actual
workflow prices and production service levels remain future
qualification requirements. A proposal can be complete without them. It
cannot honestly claim those outcomes without obtaining them.

\hypertarget{remaining-survey-and-prototype-gaps}{%
\subsection{Remaining survey and prototype
gaps}\label{remaining-survey-and-prototype-gaps}}

The public-data portfolio supports specific analyses, including the
earnings experiment and the UCI timing comparison in this volume. These
executions do not qualify every candidate dataset. Current-source web
snapshots still need exact data-file vintages for future analyses.
Several benchmark datasets need data-specific rights checks. Mortgage
histories require registration and use review; the HMDA publication
endpoint was inaccessible in this run. The Criteo license is a material
constraint for a commercial proposal. These affect individual options;
project-generated alternatives and the already local public data keep
the prototype work moving.

No inspected public financial source supplies the complete logged bank
servicing action process or independent real-bank linkage truth. That
finding is limited to the surveyed sources. It calls for transparent
semi-synthetic experiments now and a justified targeted measurement
later, rather than a claim of impossibility. Public reasoning benchmarks
also cannot establish a live agent's quality without an actual
controlled run, nor resolve the business owner's normative choice of
outcome and feasible intervention.

Finally, public data do not repair the open mathematical arguments for
persistent dependence, hidden states, graph-equivalence effect sets or
arbitrary learning procedures. Those questions remain scoped research
obligations. The proposal can present negative findings and restricted
supported regimes, compare alternatives and describe what evidence would
change the recommendation. That is a strong proposal standard without
pretending that an accepted bank product already exists.


% SOURCE docs/release/2026-10-29/source/docs/writing-repair/2026-10-29/investment.tex CHAR 2998


\section{A spending sequence with real alternatives}
The first expenditure buys a matched analyst comparison using the
constructed and public cases, with independently assessed scientific
quality and recorded full effort. It asks whether the small integration
changes decisions or reduces work enough to pay for its ongoing cost.
A bank-data pilot follows only if that evidence is favourable and the
bank supplies a defensible target, data access, assignment information
and a meaningful precision requirement. An operating trial then tests
the relevant interfaces and platform under its actual conditions.

Several observations would reverse or narrow the recommendation. Equal
quality and effort favour the established workflow. A commercial
entitlement supplying equivalent outputs at lower full cost favours
that route. Frequent lack of treatment support or reliable observation
truth favours better study design and data collection. A repeated
need for unsupported dependent inference limits the offer route rather
than justifying a broad claim. Evidence of consequential, reproducibly
avoided errors or lower full effort can justify extending the integration.

This is a conditional proposal with a selected next experiment. It
preserves the larger scientific and software programme without treating
its unfinished questions as already solved. The detailed comparisons
and operating evidence developed in the book allow a reader to inspect those
restrictions and alternatives before funding the next step.


% SOURCE docs/release/2026-10-29/source/docs/causal_banking_release_candidate_2026-10-29.tex CHAR 32921

\backmatter
\clearpage\phantomsection
\addcontentsline{toc}{chapter}{Bibliography}
\bibliographystyle{writing-repair/2026-10-29/causalplainnat-reading}
\bibliography{writing-repair/2026-10-29/references-0,writing-repair/2026-10-29/references-1,writing-repair/2026-10-29/references-2,discovery-repair/2026-10-29/refs}
\end{document}

